\section{Filtered colimits}
\label{section-directed-colimits}

\noindent
Colimits are easier to compute or describe when they
are over a filtered diagram. Here is the definition.

\begin{definition}
\label{definition-directed}
We say that a diagram $M : \mathcal{I} \to \mathcal{C}$ is {\it directed},
or {\it filtered} if the following conditions hold:
\begin{enumerate}
\item the category $\mathcal{I}$ has at least one object,
\item for every pair of objects $x, y$ of $\mathcal{I}$
there exist an object $z$ and morphisms $x \to z$,
$y \to z$, and
\item for every pair of objects $x, y$ of $\mathcal{I}$
and every pair of morphisms $a, b : x \to y$ of $\mathcal{I}$
there exists a morphism $c : y \to z$ of $\mathcal{I}$
such that $M(c \circ a) = M(c \circ b)$ as morphisms in $\mathcal{C}$.
\end{enumerate}
We say that an index category $\mathcal{I}$ is {\it directed}, or
{\it filtered} if $\text{id} : \mathcal{I} \to \mathcal{I}$ is filtered
(in other words you erase the $M$ in part (3) above).
\end{definition}

\noindent
We observe that any diagram with filtered index category is filtered,
and this is how filtered colimits usually come about. In fact, if
$M : \mathcal{I} \to \mathcal{C}$ is a filtered diagram, then we
can factor $M$ as $\mathcal{I} \to \mathcal{I}' \to \mathcal{C}$
where $\mathcal{I}'$ is a filtered index category\footnote{Namely, let
$\mathcal{I}'$ have the same objects as $\mathcal{I}$ but
where $\Mor_{\mathcal{I}'}(x, y)$ is the quotient of $\Mor_\mathcal{I}(x, y)$
by the equivalence relation which identifies
$a, b : x \to y$ if $M(a) = M(b)$.}
such that $\colim_\mathcal{I} M$ exists if and only if
$\colim_{\mathcal{I}'} M'$ exists in which case the colimits are
canonically isomorphic.

\medskip\noindent
Suppose that $M : \mathcal{I} \to \textit{Sets}$ is a filtered diagram. In
this case we may describe the equivalence relation in the formula
$$
\colim_\mathcal{I} M
=
(\coprod\nolimits_{i\in I} M_i)/\sim
$$
simply as follows
$$
m_i \sim m_{i'}
\Leftrightarrow
\exists i'', \phi : i \to i'', \phi': i' \to i'',
M(\phi)(m_i) = M(\phi')(m_{i'}).
$$
In other words, two elements are equal in the colimit if and only if
they ``eventually become equal''.

\begin{lemma}
\label{lemma-directed-commutes}
Let $\mathcal{I}$ and $\mathcal{J}$ be index categories.
Assume that $\mathcal{I}$ is filtered and $\mathcal{J}$ is finite.
Let $M : \mathcal{I} \times \mathcal{J} \to \textit{Sets}$,
$(i, j) \mapsto M_{i, j}$ be a diagram of diagrams of sets.
In this case
$$
\colim_i \lim_j M_{i, j}
=
\lim_j \colim_i M_{i, j}.
$$
In particular, colimits over $\mathcal{I}$ commute with finite products,
fibre products, and equalizers of sets.
\end{lemma}

\begin{proof}
Omitted. In fact, it is a fun exercise to prove that a category is
filtered if and only if colimits over the category commute with finite
limits (into the category of sets).
\end{proof}

\noindent
We give a counter example to the lemma in
the case where $\mathcal{J}$ is infinite. Namely, let
$\mathcal{I}$ consist of $\mathbf{N} = \{1, 2, 3, \ldots\}$
with a unique morphism $i \to i'$ whenever $i \leq i'$.
Let $\mathcal{J}$ be the discrete category
$\mathbf{N} = \{1, 2, 3, \ldots\}$ (only morphisms are identities).
Let $M_{i, j} = \{1, 2, \ldots, i\}$ with obvious inclusion maps
$M_{i, j} \to M_{i', j}$ when $i \leq i'$. In this case
$\colim_i M_{i, j} = \mathbf{N}$ and hence
$$
\lim_j \colim_i M_{i, j}
=
\prod\nolimits_j \mathbf{N}
=
\mathbf{N}^\mathbf{N}
$$
On the other hand $\lim_j M_{i, j} = \prod\nolimits_j M_{i, j}$ and
hence
$$
\colim_i \lim_j M_{i, j}
=
\bigcup\nolimits_i \{1, 2, \ldots, i\}^{\mathbf{N}}
$$
which is smaller than the other limit.

\begin{lemma}
\label{lemma-cofinal-in-filtered}
\begin{reference}
\cite[Proposition 3.2.4]{KS}
\end{reference}
Let $\mathcal{I}$ be a category. Let $\mathcal{J}$ be a full subcategory.
Assume that $\mathcal{I}$ is filtered. Assume also that for any object
$i$ of $\mathcal{I}$, there exists a morphism $i \to j$
to some object $j$ of $\mathcal{J}$. Then $\mathcal{J}$
is filtered and cofinal in $\mathcal{I}$.
\end{lemma}

\begin{proof}
Omitted. Pleasant exercise of the notions involved.
\end{proof}

\noindent
It turns out we sometimes need a more finegrained control over the
possible conditions one can impose on index categories. Thus we add
some lemmas on the possible things one can require.

\begin{lemma}
\label{lemma-preserve-products}
Let $\mathcal{I}$ be an index category, i.e., a category. Assume
that for every pair of objects $x, y$ of $\mathcal{I}$
there exist an object $z$ and morphisms $x \to z$ and $y \to z$.
Then
\begin{enumerate}
\item If $M$ and $N$ are diagrams of sets over $\mathcal{I}$,
then $\colim (M_i \times N_i) \to \colim M_i \times \colim N_i$
is surjective,
\item in general colimits of diagrams of sets over $\mathcal{I}$
do not commute with finite nonempty products.
\end{enumerate}
\end{lemma}

\begin{proof}
Proof of (1). Let $(\overline{m}, \overline{n})$
be an element of $\colim M_i \times \colim N_i$.
Then we can find $m \in M_x$ and $n \in N_y$ for some
$x, y \in \Ob(\mathcal{I})$ such that $m$ maps to
$\overline{m}$ and $n$ maps to $\overline{n}$. See
Section \ref{section-limit-sets}.
Choose $a : x \to z$ and $b : y \to z$
in $\mathcal{I}$. Then $(M(a)(m), N(b)(n))$ is an element of
$(M \times N)_z$ whose image in $\colim (M_i \times N_i)$
maps to $(\overline{m}, \overline{n})$ as desired.

\medskip\noindent
Proof of (2). Let $G$ be a non-trivial group and
let $\mathcal{I}$ be the one-object category with endomorphism monoid $G$.
Then $\mathcal{I}$ trivially satisfies the condition stated in the lemma.
Now let $G$ act on itself by translation and view the $G$-set $G$
as a set-valued $\mathcal{I}$-diagram. Then
$$
\colim_\mathcal{I} G \times \colim_\mathcal{I} G \cong G/G \times G/G
$$
is not isomorphic to
$$
\colim_\mathcal{I} (G \times G) \cong (G \times G)/G
$$
This example indicates that you cannot just drop the additional
condition Lemma \ref{lemma-directed-commutes}
even if you only care about finite products.
\end{proof}

\begin{lemma}
\label{lemma-colimits-abelian-as-sets}
Let $\mathcal{I}$ be an index category, i.e., a category. Assume
that for every pair of objects $x, y$ of $\mathcal{I}$
there exist an object $z$ and morphisms $x \to z$ and $y \to z$.
Let $M : \mathcal{I} \to \textit{Ab}$ be a diagram of abelian
groups over $\mathcal{I}$. Then the colimit of $M$ in the category
of sets surjects onto the colimit of $M$ in the category of
abelian groups.
\end{lemma}

\begin{proof}
Recall that the colimit in the category of sets is the quotient of
the disjoint union $\coprod M_i$ by relation, see
Section \ref{section-limit-sets}.
Similarly, the colimit in the category of abelian groups is a quotient
of the direct sum $\bigoplus M_i$.
The assumption of the lemma means that given $i, j \in \Ob(\mathcal{I})$
and $m \in M_i$ and $n \in M_j$, then we can find an object
$k$ and morphisms $a : i \to k$ and $b : j \to k$.
Thus $m + n$ is represented in the colimit by the element
$M(a)(m) + M(b)(n)$ of $M_k$. Thus the $\coprod M_i$
surjects onto the colimit.
\end{proof}

\begin{lemma}
\label{lemma-split-into-connected}
Let $\mathcal{I}$ be an index category, i.e., a category. Assume
that for every solid diagram
$$
\xymatrix{
x \ar[d] \ar[r] & y \ar@{..>}[d] \\
z \ar@{..>}[r] & w
}
$$
in $\mathcal{I}$ there exist an object $w$ and dotted arrows
making the diagram commute. Then $\mathcal{I}$ is either empty
or a nonempty disjoint union of connected categories having
the same property.
\end{lemma}

\begin{proof}
If $\mathcal{I}$ is the empty category, then the lemma is true.
Otherwise, we define a relation on objects of $\mathcal{I}$ by
saying that $x \sim y$ if there exist a $z$ and
morphisms $x \to z$ and $y \to z$. This is an equivalence
relation by the assumption of the lemma. Hence $\Ob(\mathcal{I})$
is a disjoint union of equivalence classes. Let $\mathcal{I}_j$
be the full subcategories corresponding to these equivalence classes.
Then $\mathcal{I} = \coprod \mathcal{I}_j$ with $\mathcal{I}_j$
nonempty as desired.
\end{proof}

\begin{lemma}
\label{lemma-preserve-injective-maps}
Let $\mathcal{I}$ be an index category, i.e., a category. Assume
that for every solid diagram
$$
\xymatrix{
x \ar[d] \ar[r] & y \ar@{..>}[d] \\
z \ar@{..>}[r] & w
}
$$
in $\mathcal{I}$ there exist an object $w$ and dotted arrows
making the diagram commute. Then
\begin{enumerate}
\item an injective morphism $M \to N$ of diagrams of sets over
$\mathcal{I}$ gives rise to an injective map $\colim M_i \to \colim N_i$
of sets,
\item in general the same is not the case for diagrams of abelian
groups and their colimits.
\end{enumerate}
\end{lemma}

\begin{proof}
If $\mathcal{I}$ is the empty category, then the lemma is true.
Thus we may assume $\mathcal{I}$ is nonempty. In this case
we can write $\mathcal{I} = \coprod \mathcal{I}_j$ where each
$\mathcal{I}_j$ is nonempty and satisfies the same property, see
Lemma \ref{lemma-split-into-connected}. Since
$\colim_\mathcal{I} M = \coprod_j \colim_{\mathcal{I}_j} M|_{\mathcal{I}_j}$
this reduces the proof of (1) to the connected case.

\medskip\noindent
Assume $\mathcal{I}$ is connected and $M \to N$ is injective, i.e.,
all the maps $M_i \to N_i$ are injective.
We identify $M_i$ with the image of $M_i \to N_i$, i.e.,
we will think of $M_i$ as a subset of $N_i$.
We will use the description of the colimits given in
Section \ref{section-limit-sets} without further mention.
Let $s, s' \in \colim M_i$ map to the same element of $\colim N_i$.
Say $s$ comes from an element $m$ of $M_i$ and $s'$ comes from an
element $m'$ of $M_{i'}$. Then we can find a sequence
$i = i_0, i_1, \ldots, i_n = i'$ of objects of $\mathcal{I}$
and morphisms
$$
\xymatrix{
&
i_1 \ar[ld] \ar[rd] & &
i_3 \ar[ld] & &
i_{2n-1} \ar[rd] & \\
i = i_0 & &
i_2 & &
\ldots & &
i_{2n} = i'
}
$$
and elements $n_{i_j} \in N_{i_j}$ mapping to each other under
the maps $N_{i_{2k-1}} \to N_{i_{2k-2}}$ and $N_{i_{2k-1}}
\to N_{i_{2k}}$ induced from the maps in $\mathcal{I}$ above
with $n_{i_0} = m$ and $n_{i_{2n}} = m'$. We will prove by induction
on $n$ that this implies $s = s'$. The base case $n = 0$ is trivial.
Assume $n \geq 1$. Using the assumption on $\mathcal{I}$
we find a commutative diagram
$$
\xymatrix{
& i_1 \ar[ld] \ar[rd] \\
i_0 \ar[rd] & & i_2 \ar[ld] \\
& w
}
$$
We conclude that $m$ and $n_{i_2}$ map to the same element of $N_w$
because both are the image of the element $n_{i_1}$.
In particular, this element is an element $m'' \in M_w$ which
gives rise to the same element as $s$ in $\colim M_i$.
Then we find the chain
$$
\xymatrix{
&
i_3 \ar[ld] \ar[rd] & &
i_5 \ar[ld] & &
i_{2n-1} \ar[rd] & \\
w & &
i_4 & &
\ldots & &
i_{2n} = i'
}
$$
and the elements $n_{i_j}$ for $j \geq 3$ which has a smaller length
than the chain we started with. This proves the induction step and the
proof of (1) is complete.

\medskip\noindent
Let $G$ be a group and let $\mathcal{I}$ be the one-object category with
endomorphism monoid $G$. Then $\mathcal{I}$ satisfies the condition stated
in the lemma because given $g_1, g_2 \in G$ we can find $h_1, h_2 \in G$
with $h_1 g_1 = h_2 g_2$. A diagram $M$ over $\mathcal{I}$ in
$\textit{Ab}$ is the same thing as an abelian group $M$ with $G$-action
and $\colim_\mathcal{I} M$ is the coinvariants $M_G$ of $M$.
Take $G$ the group of order $2$ acting trivially on $M = \mathbf{Z}/2\mathbf{Z}$
mapping into the first summand of
$N = \mathbf{Z}/2\mathbf{Z} \times \mathbf{Z}/2\mathbf{Z}$
where the nontrivial element of $G$ acts by
$(x, y) \mapsto (x + y, y)$. Then $M_G \to N_G$ is zero.
\end{proof}

\begin{lemma}
\label{lemma-split-into-directed}
Let $\mathcal{I}$ be an index category, i.e., a category.
Assume
\begin{enumerate}
\item for every pair of morphisms $a : w \to x$ and $b : w \to y$
in $\mathcal{I}$ there exist an object $z$ and morphisms $c : x \to z$
and $d : y \to z$ such that $c \circ a = d \circ b$, and
\item for every pair of morphisms $a, b : x \to y$ there exists
a morphism $c : y \to z$ such that $c \circ a = c \circ b$.
\end{enumerate}
Then $\mathcal{I}$ is a (possibly empty) union
of disjoint filtered index categories $\mathcal{I}_j$.
\end{lemma}

\begin{proof}
If $\mathcal{I}$ is the empty category, then the lemma is true.
Otherwise, we define a relation on objects of $\mathcal{I}$ by
saying that $x \sim y$ if there exist a $z$ and
morphisms $x \to z$ and $y \to z$. This is an equivalence
relation by the first assumption of the lemma. Hence $\Ob(\mathcal{I})$
is a disjoint union of equivalence classes. Let $\mathcal{I}_j$
be the full subcategories corresponding to these equivalence classes.
The rest is clear from the definitions.
\end{proof}

\begin{lemma}
\label{lemma-almost-directed-commutes-equalizers}
Let $\mathcal{I}$ be an index category satisfying the hypotheses of
Lemma \ref{lemma-split-into-directed} above. Then colimits over $\mathcal{I}$
commute with fibre products and equalizers in sets (and more generally
with finite connected limits).
\end{lemma}

\begin{proof}
By
Lemma \ref{lemma-split-into-directed}
we may write $\mathcal{I} = \coprod \mathcal{I}_j$ with each $\mathcal{I}_j$
filtered. By
Lemma \ref{lemma-directed-commutes}
we see that colimits of $\mathcal{I}_j$ commute with equalizers and
fibre products. Thus it suffices to show that equalizers and fibre products
commute with coproducts in the category of sets (including empty coproducts).
In other words, given a set $J$ and sets $A_j, B_j, C_j$ and set maps
$A_j \to B_j$, $C_j \to B_j$ for $j \in J$ we have to show that
$$
(\coprod\nolimits_{j \in J} A_j)
\times_{(\coprod\nolimits_{j \in J} B_j)}
(\coprod\nolimits_{j \in J} C_j)
=
\coprod\nolimits_{j \in J} A_j \times_{B_j} C_j
$$
and given $a_j, a'_j : A_j \to B_j$ that
$$
\text{Equalizer}(
\coprod\nolimits_{j \in J} a_j,
\coprod\nolimits_{j \in J} a'_j)
=
\coprod\nolimits_{j \in J}
\text{Equalizer}(a_j, a'_j)
$$
This is true even if $J = \emptyset$. Details omitted.
\end{proof}
